\documentclass[14pt,a4paper]{extarticle}
\usepackage[a4paper,top=12mm,bottom=12mm,left=14mm,right=14mm]{geometry}
\usepackage[T1]{fontenc}
\usepackage{amsmath}
\usepackage{amssymb}
\usepackage{xcolor}
\usepackage{tikz}
\usepackage[most]{tcolorbox}
\usetikzlibrary{arrows.meta,decorations.pathreplacing}

% ================================================================
% Algebraic terminology: a one-page sheet for a Year 7 class meeting
% algebra for the first time.
%
% Part 1 labels the parts of an expression. A labelled example of
% 5x+3 sits beside the definitions, then Q1 is the same picture of
% 7y-2 with the names left blank. Q2 labels a longer expression with
% two variables, and asks for a term as well.
%
% Part 2 fills in the gaps in sentences about an expression. Q3 gives
% most of each sentence and Q4 gives less. Q5 goes the other way, from
% the description to the expression, and Q6 is an open task.
%
% The sheet is one side of A4 in 14pt type, so that it can also be
% printed at A5 (the type is then about 10pt). The boxes and gaps are
% sized to be written in at A5. The answers are in a comment beside
% each question.
% ================================================================

\pagestyle{empty}
\setlength{\parindent}{0pt}
\setlength{\parskip}{0pt}

% The idea the sheet is built on, boxed and only faintly tinted, so the
% sheet is cheap to print.
\newtcolorbox{keyidea}{colback=orange!3,colframe=orange!35,boxrule=0.6pt,
  left=6pt,right=6pt,top=4pt,bottom=4pt}

% The section headings, between two thin rules.
\newcommand{\parttitle}[2]{%
  {\color{blue!45!black}\rule{\linewidth}{1pt}}\\[1pt]
  {\large\bfseries Part #1: #2}\\[-8pt]
  {\color{blue!45!black}\rule{\linewidth}{0.4pt}}
  \par\vspace{2pt}}

% A question number, and gaps to write in: \lgap for a letter, \gap for
% a number or a term, \wgap for a word or an expression. The strut
% keeps the lines far enough apart to write on at A5.
\newcommand{\qn}[1]{\textbf{Q#1.}\ }
\newcommand{\gapstrut}{\rule[-1.5mm]{0pt}{7.5mm}}
\newcommand{\lgap}{\gapstrut\underline{\hspace{1.1cm}}}
\newcommand{\gap}{\gapstrut\underline{\hspace{1.4cm}}}
\newcommand{\wgap}{\gapstrut\underline{\hspace{3.4cm}}}

% A Part 1 question: the number on the left, the picture centred.
\newcommand{\labelq}[2]{%
  \par\noindent\begin{minipage}[c]{1.3cm}\qn{#1}\end{minipage}\hfill
  #2\hfill\hspace*{1.3cm}\par}

% ================================================================
% Labelled expressions
% ================================================================

\tikzset{
  % The height and depth are in points, because an em here would be
  % measured in the body font and come out shorter than the letters.
  part/.style={inner xsep=0pt, inner ysep=2pt, font=\LARGE,
               text height=15pt, text depth=5pt},
  arr/.style={-{Stealth[length=7pt]}, thick, shorten >=3pt},
  brc/.style={decorate, decoration={brace, amplitude=6pt, raise=3pt}, thick},
  brcd/.style={decorate, decoration={brace, amplitude=6pt, raise=3pt,
               mirror}, thick},
}

% \anatomy[<box width>][<arrow length>]{<coefficient>}{<variable>}{<constant term>}
%         {<coefficient name>}{<variable name>}{<constant name>}{<expression name>}
% An expression of the form 5x+3 with an arrow to each part and a
% brace under the whole expression, each ending in a box. The box
% holds the name of the part, or is left empty for the name to be
% written in. The constant is written with its sign, because the sign
% in front of a term belongs to the term.
\NewDocumentCommand{\anatomy}{O{3.8cm}O{0.6cm}mmmmmmm}{%
  \begin{tikzpicture}[baseline=(v.base),
      namebox/.style={draw, minimum width=#1, minimum height=0.95cm,
                      inner sep=2pt}]
    \node[part] (c) {$#3$};
    \node[part, anchor=base west] (v) at (c.base east) {$#4$};
    \node[part, anchor=base west] (k) at ([xshift=9pt]v.base east) {$#5$};
    \node[namebox, anchor=east] (cl) at ([xshift=-#2]c.west) {#6};
    \draw[arr] (cl.east) -- (c.west);
    \node[namebox, anchor=south] (vl) at ([yshift=0.45cm]v.north) {#7};
    \draw[arr] (vl.south) -- (v.north);
    \node[namebox, anchor=west] (kl) at ([xshift=#2]k.east) {#8};
    \draw[arr] (kl.west) -- (k.east);
    \draw[brcd] (c.south west) -- (k.south east) coordinate[midway] (m);
    \node[namebox, anchor=north] at ([yshift=-0.4cm]m) {#9};
  \end{tikzpicture}%
}

\begin{document}

\begin{center}
  {\LARGE\bfseries Algebraic Terminology}
\end{center}
\vspace{-4pt}

\begin{keyidea}
\begin{minipage}[c]{0.53\linewidth}
\raggedright
A \textbf{variable} is a letter that stands for a number.

\vspace{3pt}
A \textbf{term} is a number, a variable, or numbers and variables
multiplied together. The sign in front of a term belongs to the term.

\vspace{3pt}
The \textbf{coefficient} of a term is the number that multiplies the
variables in it.

\vspace{3pt}
A \textbf{constant} is a term with no variable in it.

\vspace{3pt}
An \textbf{expression} is terms added or subtracted, with no equals sign.
\end{minipage}\hfill
\begin{minipage}[c]{0.45\linewidth}
\centering
\anatomy[2.3cm][0.4cm]{5}{x}{+3}{\small coefficient}{\small variable}{\small constant}{\small expression}
\end{minipage}
\end{keyidea}

\vspace{6pt}
% ================================================================
\parttitle{1}{Write the name of each part in its box}
% ================================================================

\vspace{4pt}
% Answers: coefficient (7), variable (y), constant (-2), expression.
\labelq{1}{\anatomy{7}{y}{-2}{}{}{}{}}

\vspace{14pt}
% Answers: coefficient (4), variable (a), term (9b), constant (-6).
\labelq{2}{%
\begin{tikzpicture}[baseline=(a.base),
    namebox/.style={draw, minimum width=3.8cm, minimum height=0.95cm, inner sep=2pt}]
  \node[part] (c) {$4$};
  \node[part, anchor=base west] (a) at (c.base east) {$a$};
  \node[part, anchor=base west] (t) at ([xshift=9pt]a.base east) {$+9b$};
  \node[part, anchor=base west] (k) at ([xshift=9pt]t.base east) {$-6$};
  \node[namebox, anchor=east] (cl) at ([xshift=-0.6cm]c.west) {};
  \draw[arr] (cl.east) -- (c.west);
  \node[namebox, anchor=north] (al) at ([yshift=-0.45cm]a.south) {};
  \draw[arr] (al.north) -- (a.south);
  \draw[brc] (t.north west) -- (t.north east) coordinate[midway] (m);
  \node[namebox, anchor=south] at ([yshift=0.5cm]m) {};
  \node[namebox, anchor=west] (kl) at ([xshift=0.6cm]k.east) {};
  \draw[arr] (kl.west) -- (k.east);
\end{tikzpicture}}

\vspace{10pt}
% ================================================================
\parttitle{2}{Fill in the gaps}
% ================================================================

\begin{minipage}[t]{0.5\linewidth}
% Answers: 3 terms; x, y and z; 4; -5.
\qn{3} \quad {\Large $4xy+2z-5$}

This expression has \gap\ terms.

Its $3$ variables are \lgap, \lgap\ and \lgap.

The coefficient of the term $xy$ is \gap.

The constant is \gap.
\end{minipage}\hfill
\begin{minipage}[t]{0.47\linewidth}
% Answers: 6m, -n and 8; 6; -1; the constant.
\qn{4} \quad {\Large $6m-n+8$}

The terms are \gap, \gap\ and \gap.

The coefficient of $m$ is \gap.

The coefficient of $n$ is \gap.

The $8$ is called the \wgap.
\end{minipage}

\vspace{12pt}
% Answer: 3p-5q+4 (the terms can be in any order).
\qn{5} An expression has the variables $p$ and $q$. The coefficient of
$p$ is $3$, the coefficient of $q$ is $-5$ and the constant is $4$.
Write down the expression. \hfill \wgap

\vspace{12pt}
% Answer: open. For example 5a+2b+ab-1. It needs four terms that are
% not like terms (so that it cannot be simplified to fewer), two
% letters for the variables, and the term -1.
\qn{6} Write an expression that has $4$ terms, $2$ variables and a
constant of $-1$.\par
\hfill\gapstrut\underline{\hspace{5.5cm}}

\end{document}
